\documentclass[11pt,a4paper,oneside]{report}
\usepackage{relatorio}
\EmpresaRelatorio{Empresa Exemplo Ltda.}
\MarcaRelatorio{Marca Exemplo}
\IdentificacaoRodape{CNPJ 00.000.000/0001-00 \textperiodcentered\ Cidade, UF}
\TituloRelatorio{Tratamento de Dados Pessoais no App Exemplo}
\SubtituloRelatorio{Relatório técnico de mapeamento, evidências e pontos de atenção à luz da Lei Geral de Proteção de Dados}
\VolumeRelatorio{Volume I \textperiodcentered\ Proteção de dados}
\TituloCurto{App Exemplo \textperiodcentered\ Dados pessoais e LGPD}
\VersaoRelatorio{Versão 1.0}
\DataRelatorio{1 de janeiro de 2027}
\AutorRelatorio{Nome do Autor}
\CargoAutor{Engenharia de Software}
\DestinatarioRelatorio{Diretoria da Empresa Exemplo Ltda. e consultoria jurídica por ela indicada}
\ClassificacaoRelatorio{Confidencial}
\title{Tratamento de Dados Pessoais no App Exemplo}
\author{Nome do Autor}
\begin{document}
\FazerCapa
\tableofcontents

\chapter{Sumário executivo}\label{cap:sumario}
\begin{destaquebox}
Síntese do capítulo em três a cinco linhas: o que foi examinado, o que se encontrou e o que se pede ao leitor.
\end{destaquebox}

Este relatório descreve, com base em evidência extraída do código-fonte, do banco de dados e de testes executados, como o aplicativo trata dados pessoais. O texto distingue o que foi \destaque{verificado} do que depende de informação externa ao código\footnote{\lgpd{art.~37}: \enquote{O controlador e o operador devem manter registro das operações de tratamento de dados pessoais que realizarem, especialmente quando baseado no legítimo interesse.}}.

\noindent
\begin{minipage}{0.235\textwidth}\indicador{12}{tabelas com dado pessoal sob controle de acesso}\end{minipage}\hfill
\begin{minipage}{0.235\textwidth}\indicador{48/50}{sondas de acesso indevido barradas}\end{minipage}\hfill
\begin{minipage}{0.235\textwidth}\indicador{7}{terceiros que recebem dados}\end{minipage}\hfill
\begin{minipage}{0.235\textwidth}\indicador{30 dias}{entre o pedido e a eliminação da conta}\end{minipage}

\section{Como ler este documento}
A gravidade é técnica: \gravidade{critica} \gravidade{alta} \gravidade{media} \gravidade{baixa} \gravidade{informativa}. A situação de cada tema aparece como \situacaotexto{atende}, \situacaotexto{parcial}, \situacaotexto{naoatende}, \situacaotexto{externo} ou \situacaotexto{na}.

\begin{itemize}
  \item A coluna \cod{accepted_at} guarda o instante do aceite; a gravação está em \arq{src/services/termos.ts:42}.
  \item Um caminho longo quebra sem estourar a margem: \arq{db/migrations/20270101120000_politica_de_retencao_de_registros.sql:106}.
  \begin{itemize}
    \item Segundo nível de lista.
  \end{itemize}
\end{itemize}

\begin{achado}{LGPD-01}{O aplicativo não apresenta aviso de privacidade ao titular}{alta}
\campo{Fato apurado}
Nenhuma tela do aplicativo nem a página pública de download informa finalidade, forma e duração do tratamento, identificação do controlador e direitos do titular. A busca por \cod{privacidade} em \arq{src/} não devolve resultado.
\campo{Evidência}
\arq{src/screens/Cadastro.tsx:1-120}; \arq{src/screens/Termo.tsx}.
\campo{Dispositivos possivelmente relacionados (hipótese)}
\lgpd{art.~6º, VI}; \lgpd{art.~9º}.
\campo{Recomendação técnica}
Publicar aviso de privacidade e exibi-lo antes do aceite, com registro da versão apresentada.
\end{achado}

\begin{pontoparecer}
A qualificação da base legal aplicável ao tratamento dos documentos enviados no cadastro depende de interpretação do contrato firmado pelo usuário.
\end{pontoparecer}

\begin{nota}
Os números de produção foram obtidos por consultas agregadas, sem leitura de registros individuais.
\end{nota}

\begin{atencao}
O ambiente de homologação contém usuários sintéticos criados por teste de carga.
\end{atencao}

\begin{conforme}
Toda tabela com dado de outro usuário é protegida por política de segurança no próprio banco.
\end{conforme}

\begin{licao}
Refutar um achado antes de classificá-lo evitou conclusões erradas.
\end{licao}

\begin{pratica}
Escreva o teste do ator não autorizado antes do caminho feliz.
\end{pratica}

\begin{pendente}
Existência de contrato de operador com cada fornecedor listado no Quadro~\ref{qd:operadores}.
\end{pendente}

\begin{quesito}{Base legal para o tratamento de dados de terceiros}
\campo{Fato que motiva a pergunta}
O resumo exibido ao usuário pode conter informação de saúde de terceiros (\achadoref{LGPD-01}).
\campo{Pergunta}
O tratamento desses dados pela empresa encontra amparo no art.~11, II, da LGPD? Em qual alínea?
\end{quesito}

\begin{table}[H]
\tabelapequena
\caption{Operadores e localização do tratamento}\label{qd:operadores}
\begin{tabularx}{\textwidth}{L{2.6cm}YL{2.4cm}C{2.2cm}}
\toprule
\linhacab \cab{Fornecedor} & \cab{Dados enviados} & \cab{Local} & \cab{Situação}\\
\midrule
Banco gerenciado & Cadastro, documentos e notificações & São Paulo (Brasil) & \situacao{atende}\\
Analytics & Identificador, e-mail, telas visitadas & Estados Unidos & \situacao{parcial}\\
Modelo de linguagem & Texto digitado no chat & Estados Unidos & \situacao{externo}\\
\bottomrule
\end{tabularx}
\end{table}

\begin{codigo}
select count(*) from public.profiles where deleted_at is not null
\end{codigo}

\section{Segunda seção}
\subsection{Subseção}
\subsubsection{Tópico}
Texto de exemplo para conferir a hierarquia tipográfica e o espaçamento entre parágrafos, com um número em destaque: \numero{242 ms}.

\part{Testes}
\chapter{Fichas}
Seta → maior ou igual ≥ e aproximadamente ≈ no texto.
\begin{ficha}{S1}{Matriz de controle de acesso com ator não autorizado}{passou}
\campo{Objetivo}
Provar que um usuário não lê nem altera dado de outro.
\campo{Resultado}
48 de 50 sondas com o resultado esperado. \resultado{falhou} \resultado{parcial} \resultado{bloqueado} \resultado{naoexecutado}
\end{ficha}
\assinatura{Nome do Autor}{Engenharia de Software}
\end{document}
